% ---------------------------------------------------------------------------
% supercm.tex -- The 3D SuperCM clustering loss: math + implementation
%
% Compile with:   pdflatex supercm.tex   (twice, for the frame numbers)
% Only standard packages: beamer, amsmath, amssymb, listings, xcolor, booktabs.
% ---------------------------------------------------------------------------
\documentclass[10pt,aspectratio=169]{beamer}

\usetheme{Madrid}
\usecolortheme{seahorse}
\setbeamertemplate{navigation symbols}{}
\setbeamertemplate{footline}[frame number]
\setbeamerfont{frametitle}{size=\large}

\usepackage[T1]{fontenc}
\usepackage{amsmath,amssymb}
\usepackage{booktabs}
\usepackage{listings}
\usepackage{xcolor}

\definecolor{codebg}{RGB}{245,245,245}
\definecolor{codekw}{RGB}{0,90,160}
\definecolor{codestr}{RGB}{160,60,0}
\definecolor{codecom}{RGB}{110,110,110}
\definecolor{hl}{RGB}{170,0,0}

\lstdefinestyle{py}{
  language=Python,
  basicstyle=\ttfamily\scriptsize,
  keywordstyle=\color{codekw}\bfseries,
  stringstyle=\color{codestr},
  commentstyle=\color{codecom}\itshape,
  backgroundcolor=\color{codebg},
  columns=flexible,
  keepspaces=true,
  showstringspaces=false,
  breaklines=true,
  frame=single,
  framesep=3pt,
  rulecolor=\color{gray},
  aboveskip=2pt,
  belowskip=2pt,
  numbers=none,
  escapeinside={(*@}{@*)},
}
\lstset{style=py}

\newcommand{\xv}{\mathbf{x}}
\newcommand{\muv}{\boldsymbol{\mu}}
\newcommand{\zv}{\mathbf{z}}
\newcommand{\Rb}{\mathbb{R}}
\newcommand{\degs}{^{\circ}}

\title[SuperCM 3D]{The SuperCM Clustering Loss in 3D}
\subtitle{From per-voxel embeddings to the contrastive CM energy\\
\small Math and implementation, side by side}
\author{MVSeg2023 / TEE few-shot segmentation}
\date{\today}

\begin{document}

\begin{frame}
  \titlepage
\end{frame}

% ===========================================================================
\begin{frame}{Notation and shapes}
\small
Everything below is for one training step of the actual configuration.

\begin{center}
\begin{tabular}{llll}
\toprule
Symbol & Meaning & Shape & Value here \\
\midrule
$B$                & batch (crops per rank)        & --                  & $2$ \\
$(H,W,D)$          & patch size                    & --                  & $112\times96\times128$ \\
$P$                & voxels per crop               & --                  & $1{,}376{,}256$ \\
$C$                & classes                       & --                  & $3$ \\
$d$                & embedding dim                 & --                  & $128$ \\
$\xv_p$            & voxel embedding               & $\Rb^{d}$           & \texttt{embedding\_dim} \\
$\zv_p$            & voxel logits                  & $\Rb^{C}$           & \\
$g_{pk}$           & softmax assignment            & $(B,C,P)$           & from $W$ \\
$y_{pk}$           & ground-truth one-hot          & $(B,C,P)$           & \texttt{cm\_mode='l'} \\
$\muv_k$           & class centroid                & $(C,d)=(3,128)$     & buffer, EMA \\
$\hat g_{pk}$      & distance-based posterior      & $(B,C,P)$           & \emph{the reconstruction} \\
$v_p$              & valid mask                    & $(B,P)$             & all ones (no ignore label) \\
\bottomrule
\end{tabular}
\end{center}

\vspace{0.4em}
Per optimizer step the loss therefore touches $B\cdot P = 2{,}752{,}512$ voxels, each a point in $\Rb^{128}$.
\end{frame}

% ===========================================================================
\section{1. The embedding}

\begin{frame}[fragile]{1. Where the 128-dimensional per-voxel feature comes from}
\begin{columns}[t]
\begin{column}{0.52\textwidth}
\footnotesize
The encoder--decoder emits \emph{stage features} at five scales. At stage $s$, voxel $p$ carries
\[
  \mathbf{f}^{(s)}_p \in \Rb^{C_s},\qquad C_s \in \{320,256,128,64,32\}.
\]
Each stage's \texttt{seg\_layer} is a $1{\times}1{\times}1$ convolution, i.e.\ a \emph{per-voxel linear map}:
\[
  \boxed{\;\xv^{(s)}_p = W^{(s)}_{\mathrm{seg}}\,\mathbf{f}^{(s)}_p + \mathbf{b}^{(s)}_{\mathrm{seg}}\;}
\]
with $W^{(s)}_{\mathrm{seg}} \in \Rb^{128 \times C_s}$.

\vspace{0.4em}
Stock nnU-Net sets the output width to $C=3$. We pass \texttt{embedding\_dim}$=128$ in its place --- that single
substituted argument is the whole architectural change.

\vspace{0.4em}
The CM loss uses \textbf{stage $0$ only} (full resolution), so from here on
$\xv_p \equiv \xv^{(0)}_p \in \Rb^{128}$.
\end{column}

\begin{column}{0.46\textwidth}
\begin{lstlisting}
# module/nnunet.py -- build_network
backbone = nnUNetTrainer.build_network_architecture(
    self.pm, self.cm,
    self.num_input_channels,
    embedding_dim,          # 128, not num_classes
    self.enable_deep_supervision,
)
network = EmbeddingAdapter(
    backbone,
    embedding_dim=embedding_dim,   # 128
    num_classes=self.num_classes,  # 3
    conv_op=conv_op,
)
\end{lstlisting}
\begin{lstlisting}
# module/adapter.py
def forward(self, x):
    embeddings = self.backbone(x)   # list, 5 stages
    logits = self._classify(embeddings)
    if self.return_embeddings:
        return embeddings, logits
    return logits
\end{lstlisting}
\footnotesize
Stage 0 embedding: \texttt{(2, 128, 112, 96, 128)}.
\end{column}
\end{columns}
\end{frame}

% ===========================================================================
\section{2. Assignments}

\begin{frame}[fragile]{2. Logits and the softmax assignment $g$}
\begin{columns}[t]
\begin{column}{0.52\textwidth}
\footnotesize
One \textbf{shared} classifier $W \in \Rb^{3\times128}$, $\mathbf{b}\in\Rb^{3}$, applied at every stage:
\[
  \zv_p = W\xv_p + \mathbf{b} \in \Rb^{3},
\]
\[
  \boxed{\;g_{pk} = \frac{\exp(z_{pk})}{\sum_{j=1}^{3}\exp(z_{pj})}\;}
\]

\vspace{0.3em}
Ground-truth one-hot targets from the integer label map $\ell_p$:
\[
  y_{pk} = \mathbf{1}[\ell_p = k]\,v_p,
  \qquad
  v_p = \mathbf{1}[\ell_p \neq \text{ignore}].
\]

\vspace{0.5em}
\textcolor{hl}{Key observation.} $g$ depends on the embedding \emph{only through} $W$. The clustering
geometry and the decision function are therefore two separate objects --- this is the seed of every
failure mode on slide~\ref{frame:collapse}.
\end{column}

\begin{column}{0.46\textwidth}
\begin{lstlisting}
@staticmethod
def _flatten_spatial(tensor):
    """(B, C, *spatial) -> (B, C, P)."""
    return tensor.flatten(2)

def _probabilities(self, logits):
    # float32 regardless of autocast: these feed
    # a reduction over ~1e6 voxels, where fp16
    # summation drifts.
    return torch.softmax(
        self._flatten_spatial(logits).float(),
        dim=1,
    )

def _gt_assignments(self, target):
    labels = target[:, 0].reshape(
        target.shape[0], -1).long()
    valid = self._valid_mask(labels)
    safe = labels.clamp(0, self.num_classes - 1)
    y = nn.functional.one_hot(
        safe, num_classes=self.num_classes)
    y = y.permute(0, 2, 1).to(torch.float32)
    return y * valid.unsqueeze(1), valid
\end{lstlisting}
\end{column}
\end{columns}
\end{frame}

% ===========================================================================
\section{3. Centroids}

\begin{frame}[fragile]{3. Centroid estimation --- the first einsum}
\begin{columns}[t]
\begin{column}{0.50\textwidth}
\footnotesize
The assignment-weighted mean, summed over all DDP ranks $r$ \emph{before} dividing:
\[
  N_k = \sum_{r}\sum_{p} y^{(r)}_{pk}\,\xv^{(r)}_p \in \Rb^{128},
  \qquad
  D_k = \sum_{r}\sum_{p} y^{(r)}_{pk}
\]
\[
  \boxed{\;\hat\muv_k = \frac{N_k}{\max(D_k,\;\varepsilon)}\;}
\]

The numerator is exactly the contraction
\[
  N_{kd} \;=\; \sum_{b}\sum_{p} x_{bdp}\, y_{bkp}
  \;\;\equiv\;\;
  \texttt{"bdp,bcp->cd"} .
\]

\vspace{0.4em}
Cost: $B\,d\,C\,P \approx 2$ GFLOP, and under \texttt{no\_grad} it keeps \emph{no} activations ---
the entire centroid estimate is one einsum with zero memory footprint.
\end{column}

\begin{column}{0.48\textwidth}
\begin{lstlisting}
@torch.no_grad()
def update_centroids(self, streams):
    numerator = torch.zeros_like(self.mu)
    denominator = torch.zeros(
        self.num_classes,
        device=self.mu.device,
        dtype=self.mu.dtype)

    for x, y in streams:
        # autocast OFF, fp32 forced. torch.einsum is
        # on autocast's fp16 list, so an explicit
        # .float() alone is SILENTLY OVERRIDDEN.
        with torch.autocast(
                device_type=x.device.type,
                enabled=False):
            x32 = x.float()
            y32 = y.float()
            numerator += torch.einsum(
                (*@\textcolor{hl}{\texttt{"bdp,bcp->cd"}}@*), x32, y32)
            denominator += y32.sum(dim=(0, 2))

    if torch.distributed.is_initialized():
        torch.distributed.all_reduce(numerator)
        torch.distributed.all_reduce(denominator)
\end{lstlisting}
\end{column}
\end{columns}
\end{frame}

% ===========================================================================
\begin{frame}[fragile]{3b. Why fp32 and why all-reduce before the division}
\begin{columns}[t]
\begin{column}{0.50\textwidth}
\footnotesize
\textbf{fp16 overflow.} The reduction runs over $2{,}752{,}512$ voxels; fp16 saturates at $65504$, so the
running total hits $\infty$ almost immediately. Then
\[
  \muv_k = \infty
  \;\Longrightarrow\;
  \frac{\partial E}{\partial g} \ni \lVert\muv_k\rVert^2 = \infty
  \;\Longrightarrow\;
  0 \cdot \infty = \mathrm{NaN}
\]
in the backward pass --- in 102 of 110 parameter tensors, \emph{even at} $\lambda=0$, because
$0.0\cdot\infty$ is also NaN. AMP's \texttt{GradScaler} then skips every optimizer step and nothing trains.

\vspace{0.5em}
\textbf{All-reduce placement.} Reducing $N_k$ and $D_k$ separately and dividing afterwards gives every rank the
\emph{same} $\muv$. Dividing per rank first would make the loss a different function on each rank: gradients
would still synchronise, but they would be gradients of eight different objectives.

\vspace{0.5em}
\textbf{EMA}, with $\alpha = 0.95$:
\[
  \muv_k \leftarrow
  \begin{cases}
    \hat\muv_k, & \text{first batch containing class } k,\\[2pt]
    \alpha\muv_k + (1-\alpha)\hat\muv_k, & \text{thereafter.}
  \end{cases}
\]
A cold EMA from $\mathbf{0}$ would need ${\sim}100$ steps to approach the true mean.
\end{column}

\begin{column}{0.48\textwidth}
\begin{lstlisting}
present = denominator > self.eps
if not bool(present.any()):
    return

estimate = (numerator[present]
            / denominator[present].unsqueeze(1))

# Fail loudly rather than poison the run: a
# non-finite centroid clips the forward energy to
# LOSS_CLIP_VALUE (looks merely large in the logs)
# while the backward carries 0 * inf = NaN into
# every parameter.
if not torch.isfinite(estimate).all():
    raise RuntimeError(
        "ClusteringCMLoss: non-finite centroid "
        "estimate. ...")

rows = present.nonzero(as_tuple=True)[0]
fresh = ~self.mu_initialized[rows]
blended = torch.where(
    fresh.unsqueeze(1),
    estimate,
    self.ema_weight * self.mu[rows]
    + (1.0 - self.ema_weight) * estimate,
)
\end{lstlisting}
\end{column}
\end{columns}
\end{frame}

% ===========================================================================
\section{4. Distances and $\hat g$}

\begin{frame}[fragile]{4. Distances and the reconstructed posterior $\hat g$ --- the second einsum}
\begin{columns}[t]
\begin{column}{0.50\textwidth}
\footnotesize
With \texttt{cm\_normalize=true}, project both onto the unit sphere:
\[
  \tilde\xv_p = \frac{\xv_p}{\lVert\xv_p\rVert_2},
  \qquad
  \tilde\muv_k = \frac{\muv_k}{\lVert\muv_k\rVert_2}.
\]
Then
\[
  d^2_{pk}
  = \lVert\tilde\xv_p - \tilde\muv_k\rVert_2^2
  = \underbrace{\lVert\tilde\xv_p\rVert^2}_{=1}
    - 2\,\tilde\xv_p^{\!\top}\tilde\muv_k
    + \underbrace{\lVert\tilde\muv_k\rVert^2}_{=1}
\]
\[
  \boxed{\;d^2_{pk} = 2 - 2\cos\angle(\xv_p,\muv_k)\;}
\]
The code evaluates the \emph{expanded} middle form, never the explicit difference: the largest retained
tensor is then $(B,C,P)$ rather than $(B,d,P)$ --- a $128/3 \approx 43\times$ saving. The cross term is
\[
  \texttt{"bdp,cd->bcp"} :\quad
  \sum_{d} x_{bdp}\,\mu_{cd}.
\]

\vspace{0.3em}
The \textbf{reconstructed posterior} --- nearest-centroid assignment, built from geometry alone:
\[
  \boxed{\;\hat g_{pk} = \frac{\exp(-d^2_{pk}/\tau)}{\sum_{j}\exp(-d^2_{pj}/\tau)}\;}
\]
\end{column}

\begin{column}{0.48\textwidth}
\begin{lstlisting}
def _energy(self, x, g, y, valid):
    with torch.autocast(
            device_type=x.device.type,
            enabled=False):
        x = x.float()
        g = g.float()
        y = y.float()
        mu = self.mu.float()

        if self.normalize:
            x = nn.functional.normalize(x, dim=1)
            mu = nn.functional.normalize(mu, dim=1)

        x_squared = x.pow(2).sum(dim=1).unsqueeze(1)
        x_dot_mu = torch.einsum(
            (*@\textcolor{hl}{\texttt{"bdp,cd->bcp"}}@*), x, mu)
        mu_squared = mu.pow(2).sum(dim=1).view(1,-1,1)

        # (B, C, P) squared distance, every voxel
        # to every centroid
        dist_sq = (x_squared
                   - 2.0 * x_dot_mu
                   + mu_squared).clamp_min(0.0)
\end{lstlisting}
\footnotesize
$\tau$ is \texttt{cm\_tau}$=1.0$. Note $\hat g$ never sees $W$, and $g$ never sees $\muv$.
\end{column}
\end{columns}
\end{frame}

% ===========================================================================
\section{5. The energy}

\begin{frame}[fragile]{5. The CM energy: two forms}
\begin{columns}[t]
\begin{column}{0.50\textwidth}
\footnotesize
\textbf{Contrastive} (\texttt{cm\_form=contrastive}, current) --- cross-entropy between the target and the
\emph{distance} posterior:
\[
  \boxed{\;
  E_{\mathrm{con}}
  = \frac{\sum_p v_p\Bigl(-\sum_k y_{pk}\log\hat g_{pk}\Bigr)}{\sum_p v_p}
  \;}
\]
bounded: $0 \le E_{\mathrm{con}} \le \log 3 \approx 1.0986$.

\vspace{0.8em}
\textbf{Attract} (\texttt{cm\_form=attract}, the original SuperCM soft $k$-means energy) --- $g$-weighted
squared distance:
\[
  \boxed{\;
  E_{\mathrm{att}}
  = \frac{\sum_p v_p \sum_k g_{pk}\,d^2_{pk}}{\sum_p v_p}
  \;}
\]
unbounded above; divided by $d$ when \texttt{normalize} is off, purely to keep the scale comparable.

\vspace{0.6em}
Both are means over \emph{valid} voxels, which is why the denominator is $\sum_p v_p$ and not $BP$.
\end{column}

\begin{column}{0.48\textwidth}
\begin{lstlisting}
        if self.form == "contrastive":
            log_posterior = nn.functional.log_softmax(
                -dist_sq / self.tau, dim=1)
            per_voxel = -(y * log_posterior).sum(dim=1)
        else:
            per_voxel = (g * dist_sq).sum(dim=1)
            if not self.normalize:
                per_voxel = (per_voxel
                             / float(self.embedding_dim))

        weights = valid.to(per_voxel.dtype)
        return ((per_voxel * weights).sum()
                / weights.sum().clamp_min(1.0))
\end{lstlisting}

\vspace{0.3em}
\begin{lstlisting}
# forward(): labeled term always, unlabeled when present
energy_terms = [(x_l, g_l, y_l, valid_l)]
if x_u is not None:
    ones = torch.ones_like(g_u[:, 0], dtype=torch.bool)
    y_u = y_u if y_u is not None else g_u
    energy_terms.append((x_u, g_u, y_u, ones))

for x, g, y, valid in energy_terms:
    term = self._energy(x, g, y, valid)
    total = term if total is None else total + term
total = _clip_loss(total / float(len(energy_terms)))
\end{lstlisting}
\end{column}
\end{columns}
\end{frame}

% ===========================================================================
\begin{frame}{6. The full training objective}
\small
\[
  \mathcal{L}
  = \underbrace{\sum_{s=0}^{4} w_s\,
      \mathcal{L}_{\mathrm{Dice{+}CE}}\!\bigl(\zv^{(s)},\,\ell^{(s)}\bigr)}
    _{\text{supervised, all five deep-supervision stages}}
  \;+\;
  \underbrace{\lambda(t)\,E}_{\text{CM, stage 0 only}}
\]

\vspace{0.6em}
Deep-supervision weights, $w_s \propto 2^{-s}$ and normalised to sum to one:
\[
  w \propto \bigl[1,\tfrac12,\tfrac14,\tfrac18,10^{-6}\bigr]
  \;\Longrightarrow\;
  w = [0.533,\;0.267,\;0.133,\;0.067,\;{\approx}5\!\times\!10^{-7}].
\]
The last stage carries $10^{-6}$ rather than $0$ under DDP: a weight of exactly zero leaves the deepest
\texttt{seg\_layer} without a gradient, and DDP aborts on an unused parameter.

\vspace{0.8em}
The CM weight is a \emph{delayed constant}, not a ramp:
\[
  \lambda(t) =
  \begin{cases}
    0, & t < 5,\\
    \lambda_0 = 1.0, & t \ge 5,
  \end{cases}
  \qquad
  t = \frac{\texttt{global\_step}}{\texttt{num\_training\_batches}} .
\]
The five-epoch delay exists because $\muv$ is meaningless until the embedding has structure --- it is a
warm-up for the \emph{centroids}, not a curriculum for the loss.
\end{frame}

% ===========================================================================
\section{6. Why the form matters}

\begin{frame}{7. Take the gradient --- this is the whole story}
\label{frame:collapse}
\small
Let $u_{pk} = -d^2_{pk}/\tau$, so $\hat g_p = \mathrm{softmax}(u_p)$. For a softmax cross-entropy,
\[
  \frac{\partial E_{\mathrm{con}}}{\partial u_{pk}} = \hat g_{pk} - y_{pk},
  \qquad
  \frac{\partial u_{pk}}{\partial \tilde\xv_p} = -\frac{2}{\tau}\bigl(\tilde\xv_p - \tilde\muv_k\bigr).
\]
Chaining, the \textbf{descent} direction is
\[
  \boxed{\;
  -\frac{\partial E_{\mathrm{con}}}{\partial \tilde\xv_p}
  = -\frac{2}{\tau}\sum_k \bigl(y_{pk} - \hat g_{pk}\bigr)\bigl(\tilde\xv_p - \tilde\muv_k\bigr)
  \;}
\]

\vspace{0.4em}
Read the two cases off the sign of the residual $(y_{pk}-\hat g_{pk})$:

\begin{center}
\begin{tabular}{lll}
\toprule
$k$ & residual & effect on $\tilde\xv_p$ \\
\midrule
true class $k^\star$ & $1 - \hat g_{pk^\star} > 0$ & moves \textbf{toward} $\tilde\muv_{k^\star}$ --- attraction \\
any $k \neq k^\star$ & $-\hat g_{pk} < 0$          & moves \textbf{away from} $\tilde\muv_{k}$ --- repulsion \\
\bottomrule
\end{tabular}
\end{center}

\vspace{0.4em}
Repulsion appears \emph{automatically}, weighted by exactly how much probability mass the wrong centroid
currently claims. No extra term, no extra hyper-parameter.
\end{frame}

% ===========================================================================
\begin{frame}{8. The same calculation for the attract form}
\small
\[
  E_{\mathrm{att}} = \sum_k g_{pk}\, d^2_{pk}
  \qquad\Longrightarrow\qquad
  -\frac{\partial E_{\mathrm{att}}}{\partial \tilde\xv_p}
  = -2\sum_k g_{pk}\bigl(\tilde\xv_p - \tilde\muv_k\bigr)
\]

\vspace{0.3em}
Every coefficient $g_{pk} \ge 0$. \textcolor{hl}{\textbf{Every term is attraction.}} There is no $k$ for which
this pushes the embedding away from a centroid.

\vspace{0.6em}
Measured energy as the three centroids are rotated together (synthetic, unit-norm, $\tau=1$):

\begin{center}
\begin{tabular}{lcccccc}
\toprule
pairwise centroid angle & $120\degs$ & $90\degs$ & $60\degs$ & $20\degs$ & $5\degs$ & $1\degs$ \\
\midrule
$E_{\mathrm{att}}$  & 0.6719 & 0.6697 & 0.6650 & 0.6618 & 0.6625 & 0.6629 \\
$E_{\mathrm{con}}$  & \textbf{0.2475} & 0.3432 & 0.5794 & 0.9992 & 1.0890 & \textbf{1.0955} \\
\bottomrule
\end{tabular}
\end{center}

\vspace{0.5em}
$E_{\mathrm{att}}$ is \emph{flat} --- and in fact slightly \emph{minimised} near $20\degs$. Collapse is a
descent direction. $E_{\mathrm{con}}$ is monotone in the separation, so collapse is an ascent direction.

\vspace{0.5em}
\textbf{What was actually observed} in sweep \texttt{7yu5e5o5} with the attract form: mean pairwise centroid
angle fell $113\degs \to 12\degs$, within-class spread $25\degs \to 1\degs$, Dice $0.002$--$0.021$ \emph{worse}
than the supervised baseline.
\end{frame}

% ===========================================================================
\begin{frame}{9. Why labels do not prevent centroid collapse}
\small
The objection is natural: \emph{$\muv$ is computed from ground-truth labels, so how can it collapse?}

\vspace{0.6em}
Look again at the estimator:
\[
  \muv_k = \frac{\sum_p y_{pk}\,\xv_p}{\sum_p y_{pk}}.
\]
$\muv_k$ is a \textbf{function of} $\xv$. The labels $y$ decide \emph{which} embeddings get averaged
together; they say nothing about \emph{where} those averages sit in $\Rb^{128}$. Gradient descent is free to
move all three groups to the same place --- the partition stays correct and the energy goes down.

\vspace{0.8em}
Two failure modes followed from exactly this, both now fixed:

\begin{enumerate}
\item \textbf{Scale degeneracy.} $\xv \mapsto \xv/s$, $W \mapsto sW$ leaves every logit \emph{identical}.
      $E_{\mathrm{att}}$ falls as $1/s^2$. Observed: embedding norm shrank $1.9\times$ while $\lVert W\rVert$
      grew $1.37\times$. \;Fix: \texttt{cm\_normalize=true}, which quotients out the scale.
\item \textbf{Angular collapse.} Once scale is fixed, the remaining gauge freedom is rotation; the attract
      energy is flat in the separation. \;Fix: \texttt{cm\_form=contrastive}, which makes separation strictly
      preferred.
\end{enumerate}

\vspace{0.6em}
The 2D SuperCM project was immune to both because its head was
\texttt{pred = torch.mean(pre\_heatmaps, dim=1)} --- \emph{parameterless}. Introducing a learnable $W$ is what
opened the gauge.
\end{frame}

% ===========================================================================
\begin{frame}[fragile]{10. One step, end to end}
\begin{columns}[t]
\begin{column}{0.46\textwidth}
\footnotesize
\begin{enumerate}\itemsep0.25em
\item Forward once with embeddings on $\rightarrow$ 5 stage embeddings $+$ 5 stage logit maps.
\item Supervised loss over \emph{all} stages, via \texttt{DeepSupervisionWrapper}.
\item Stage-0 embedding and logits only $\rightarrow$ CM loss.
\item Inside CM: $g$ from logits, $y$ from labels, $\muv$ re-estimated and EMA'd (\texttt{no\_grad}),
      $d^2$ and $\hat g$ built, energy reduced.
\item $\mathcal{L} = \mathcal{L}_{\mathrm{sup}} + \lambda(t)\,E$.
\end{enumerate}

\vspace{0.5em}
The centroid update sits \emph{inside} the loss forward, under \texttt{no\_grad}, so $\muv$ is a constant with
respect to the backward pass --- the gradient flows only through $\xv_p$ in $d^2_{pk}$, exactly as the
derivation on slide~\ref{frame:collapse} assumes.
\end{column}

\begin{column}{0.52\textwidth}
\begin{lstlisting}
# module/lightning_module.py -- training_step
embeddings, logits = self.forward_embeddings(batch["data"])

supervised_loss = self.loss(logits, target)       # all 5 stages

cm_energy, cm_stats = self.cm_loss(
    self._highest_resolution(embeddings),         # stage 0
    self._highest_resolution(logits),             # stage 0
    self._highest_resolution(target),
)

cm_weight = self.cm_loss.weight_at(self._cm_epoch())
total_loss = supervised_loss + cm_weight * cm_energy
\end{lstlisting}

\begin{lstlisting}
def weight_at(self, epoch: float) -> float:
    return (self.weight
            if float(epoch) >= self.start_epoch
            else 0.0)

def _cm_epoch(self):
    n = self.trainer.num_training_batches
    if n and math.isfinite(n):
        return self.global_step / float(n)
    return float(self.current_epoch)
\end{lstlisting}
\end{column}
\end{columns}
\end{frame}

% ===========================================================================
\begin{frame}{11. Configuration map}
\small
\begin{center}
\begin{tabular}{llp{6.2cm}}
\toprule
Config key & Value & Role in the math \\
\midrule
\texttt{embedding\_dim}   & 128              & $d$; width of the \texttt{seg\_layer} output \\
\texttt{cm\_mode}         & \texttt{'l'}     & $y_{pk}$ from ground-truth labels \\
\texttt{cm\_form}         & \texttt{contrastive} & $E_{\mathrm{con}}$ instead of $E_{\mathrm{att}}$ \\
\texttt{cm\_tau}          & 1.0              & $\tau$ in $\hat g = \mathrm{softmax}(-d^2/\tau)$ \\
\texttt{cm\_weight}       & 1.0              & $\lambda_0$ \\
\texttt{cm\_start\_epoch} & 5                & the step in $\lambda(t)$ \\
\texttt{cm\_normalize}    & \texttt{true}    & $\tilde\xv, \tilde\muv$ on the unit sphere \\
\texttt{cm\_ema\_weight}  & 0.95             & $\alpha$ in the centroid EMA \\
\texttt{patch\_size\_override} & [112, 96, 128] & fixes $P = 1{,}376{,}256$ \\
\texttt{precision}        & \texttt{16-mixed} & why every CM block disables autocast explicitly \\
\bottomrule
\end{tabular}
\end{center}

\vspace{0.6em}
\footnotesize
Three keys exist purely for the unlabeled branch and are inert while it is off:
\texttt{cm\_conf\_thr}$=0.90$, \texttt{cm\_topk\_frac}$=0.15$, \texttt{cm\_min\_keep}$=64$,
\texttt{cm\_use\_soft\_pseudo}$=$\texttt{false}.
\end{frame}

% ===========================================================================
\begin{frame}{12. Check yourself}
\small
At \texttt{cm\_mode='l'}, the energy trains $\hat g$ --- the \emph{distance} posterior --- toward the
ground-truth $y$. But the network's actual prediction is
\[
  g = \mathrm{softmax}(W\xv + \mathbf{b}),
\]
and the CM loss never touches $W$ or $\mathbf{b}$.

\vspace{1.2em}
\begin{center}
\textbf{Through what path, then, does minimising $E$ change the segmentation at all?}
\end{center}

\vspace{1.2em}
There is a path. Naming it tells you precisely what this regulariser can be expected to do --- and what it
cannot.

\vspace{1.5em}
\footnotesize
\textcolor{gray}{Follow-up worth deriving: what does $\tau \to 0$ do to the gradient on slide
\ref{frame:collapse}, and why is $\tau = 1$ with unit-norm vectors ($d^2 \in [0,4]$) a reasonable default?}
\end{frame}

\end{document}
